% !TEX program = xelatex
%% ══════════════════════════════════════════════════════════════════
%%  NovaForge — 通用知识笔记 LaTeX 模板
%%  适用于：考研/考公/专业课/科研/项目/竞赛
%%  编译方式：xelatex × 2 遍后清理 .aux/.log/.out/.toc 等辅助文件
%%  使用 scripts/compile.sh（macOS/Linux）或 scripts/compile.bat（Windows）一键完成
%% ══════════════════════════════════════════════════════════════════
\documentclass[10pt,a4paper]{article}

%% ── 中文 ──
\usepackage{xeCJK}
\setCJKmainfont{SimSun}
\usepackage{indentfirst}

%% ── 中文引号（重要） ──
%% 必须使用 Unicode 弯引号：“”(U+201C/U+201D)，不可使用 ASCII 直引号 " (U+0022)
%% 直引号在 xelatex 中渲染方向错误，编辑器中请设置自动替换直引号为弯引号

%% ── 数学/公式（非理工科可移除） ──
\usepackage{amsmath,amssymb,amsthm}
\usepackage{cancel}       % \cancel{} 划掉项
\usepackage{physics}      % 向量/微分/括号简写
\usepackage{mathrsfs}     % \mathscr 花体
\usepackage{bm}           % \bm{} 粗体符号

%% ── 页面 ──
\usepackage[top=1.4cm,bottom=1.2cm,left=1.0cm,right=1.0cm,includehead]{geometry}
\usepackage{xcolor}
\usepackage{enumitem}
\usepackage{array,booktabs,multirow}
\usepackage{fancyhdr}
\usepackage{titlesec}
\usepackage{hyperref}
\usepackage{environ}
\usepackage{tikz}
\usetikzlibrary{decorations.pathmorphing,patterns,arrows.meta}

\hypersetup{hidelinks}

%% ── 颜色系统（可按学科自定义） ──
%%  理科风（深蓝） | 文科风（暗红） | 清新风（墨绿）
\definecolor{titlecolor}{HTML}{1a237e}    % 主标题 ← 从这里换色
\definecolor{seccolor}{HTML}{004d40}      % 节标题
\definecolor{subseccolor}{HTML}{b71c1c}   % 小节标题
\definecolor{emphcolor}{HTML}{bf360c}     % 强调
\definecolor{infocolor}{HTML}{1565c0}     % 提示框
\definecolor{warncolor}{HTML}{e65100}     % 警告
\definecolor{supercolor}{HTML}{880e4f}    % 超纲/拓展
\definecolor{examplecolor}{HTML}{1b5e20}  % 示例
\definecolor{sectionbg}{HTML}{e8eaf6}     % 标题背景
\definecolor{lightgray}{HTML}{f5f5f5}
\definecolor{hwcolor}{HTML}{795548}       % 习题
\definecolor{learncolor}{HTML}{2e7d32}    % 练习

%% ── 标题格式 ──
\titleformat{\section}{\LARGE\bfseries\color{titlecolor}}{}{0em}{}[\vspace{-0.2em}\rule{\textwidth}{0.8pt}]
\titleformat{\subsection}{\large\bfseries\color{seccolor}}{}{0em}{}
\titleformat{\subsubsection}{\normalsize\bfseries\color{subseccolor}}{}{0em}{}
\titlespacing*{\section}{0pt}{1em}{0.5em}
\titlespacing*{\subsection}{0pt}{0.6em}{0.2em}
\titlespacing*{\subsubsection}{0pt}{0.4em}{0.1em}

%% ── 页眉 ──
\pagestyle{fancy}
\fancyhf{}
\fancyhead[L]{\small\color{gray}<项目/考试说明>}
\fancyhead[R]{\small\color{gray}<模块名> · <作者>}
\fancyfoot[C]{\small\color{gray}\thepage}
\renewcommand{\headrulewidth}{0.4pt}
\setlength{\headheight}{14pt}

%% ── 自定义命令 ──
\newcommand{\key}[1]{\textcolor{emphcolor}{\textbf{#1}}}       % 强调
\newcommand{\super}[1]{\textcolor{supercolor}{\textbf{#1}}}    % 拓展
\newcommand{\source}[1]{\textcolor{purple!60!black}{\small [#1]}} % 标注来源
\newcommand{\formula}[1]{\vspace{0.3em}\begin{center}
  \fcolorbox{titlecolor}{white}{\parbox{0.92\textwidth}
  {\centering\color{titlecolor}\small #1}}\end{center}\vspace{0.2em}} % 原理/公式框
\newcommand{\infobox}[1]{\vspace{0.3em}{\small\color{infocolor}
  \noindent$\blacktriangleright$ #1}\vspace{0.2em}}            % 提示
\newcommand{\warning}[1]{\vspace{0.3em}{\small\color{warncolor}
  \noindent\textbf{!} #1}\vspace{0.2em}}                       % 注意
\newcommand{\sep}{\vspace{0.5em}\hrule\vspace{0.5em}}          % 分隔

%% ── 知识标题栏 ──
\newcommand{\knowtitle}[1]{%
  \vspace{0.4em}%
  \noindent\colorbox{sectionbg}{\parbox{\dimexpr\textwidth-2\fboxsep\relax}
  {\small\bfseries\color{titlecolor}#1}}%
  \vspace{0.2em}%
}

%% ── 例题/案例环境（紫色标签，自动带"解/"分析"） ──
\newenvironment{exampenv}[1]{%
  \vspace{0.5em}%
  \noindent{\color{supercolor!15}\rule{0.4em}{0pt}}\hspace{-0.4em}%
  \textcolor{supercolor}{\rule{0.4em}{0.9em}}\hspace{0.3em}%
  \textcolor{supercolor}{\textbf{【案例/例题】#1}}\par
  \noindent\hspace{0.7em}{\small\color{infocolor}分析：}\normalsize
}{%
  \par\vspace{0.3em}%
}

%% ── 练习/复盘环境（绿色标签，带解答） ──
\newenvironment{pracenv}[1]{%
  \vspace{0.5em}%
  \noindent{\color{learncolor!15}\rule{0.4em}{0pt}}\hspace{-0.4em}%
  \textcolor{learncolor}{\rule{0.4em}{0.9em}}\hspace{0.3em}%
  \textcolor{learncolor}{\textbf{【练习/复盘】#1}}\par
  \noindent\hspace{0.7em}{\small\color{gray}解答：}\normalsize
}{%
  \par\vspace{0.3em}%
}

%% ── 课后作业环境（棕色标签，解答见附录） ──
\newenvironment{hwenv}[2]{%
  \vspace{0.5em}%
  \noindent{\color{hwcolor!15}\rule{0.4em}{0pt}}\hspace{-0.4em}%
  \textcolor{hwcolor}{\rule{0.4em}{0.9em}}\hspace{0.3em}%
  \textcolor{hwcolor}{\textbf{【作业】#1}}\par
  \noindent\hspace{0.7em}{\small\color{gray}（解答$\to$ \hyperref[#2]{附录\ref*{#2}}）}\normalsize
}{%
  \par\vspace{0.3em}%
}

%% ════════════════════════════════════════════
\begin{document}

%% ── 封面 ──
\begin{center}
  \vspace{2em}
  {\Huge\bfseries\color{titlecolor} <科目/项目名称>} \\[0.3em]
  {\LARGE\bfseries\color{titlecolor} <副标题/考试全称>} \\[0.3em]
  {\large\color{gray!70!black} <一句话定位>} \\[0.3em]
  \vspace{0.3em}
  {\small\color{gray} 参考：<教材/资料1> \;$\vert$\; <教材/资料2>} \\[0.3em]
  {\small\color{gray} <作者> \quad <机构/方向> \;$\vert$\; <模块名>} \\[0.6em]
  {\normalsize\color{titlecolor!80!black}\textbf{版本：<版本号/年份>}} \\[0.1em]
  {\footnotesize\color{gray} 最后修订：<日期>}
  \vspace{1em}
\end{center}

%% ── 范围说明（如有多个方向/模块） ──
\section*{范围说明}
\noindent <说明>

\begin{center}
{\footnotesize
\begin{tabular}{@{}ll@{}}
\toprule
\textbf{模块} & \textbf{覆盖内容} \\
\midrule
<模块A> & <内容> \\
\textbf{<模块B>} & \textbf{<内容> $\leftarrow$ 本资料对应} \\
\bottomrule
\end{tabular}
}
\end{center}

%% ── 使用说明 ──
\section*{如何使用这份笔记}
\noindent 这份笔记的目标是\textbf{帮你建立系统的知识体系}：
\begin{itemize}
  \item \textbf{搭配学习}：配合教材/课程使用，先理解再巩固
  \item \textbf{独立使用}：完整覆盖核心知识点，独立查阅
  \item \textbf{复习冲刺}：方法总结 + 典型例题 + 速查表，快速回顾
\end{itemize}

\noindent\textcolor{emphcolor}{\textbf{每节结构：}}
\begin{itemize}[leftmargin=2em]
  \item[\textbf{1}] \textbf{概念引入} — 先建立直觉
  \item[\textbf{2}] \textbf{核心原理} — 公式/原理框突出
  \item[\textbf{3}] \textbf{方法技巧} — 实操技巧浓缩
  \item[\textbf{4}] \textbf{典型示例} — 经典案例搭桥
  \item[\textbf{5}] \textbf{真题/实战} — 真实题目/项目
  \item[\textbf{6}] \textbf{巩固练习} — 复盘训练
  \item[\textbf{7}] \textbf{专题总结} — 末章归纳核心题型/模式
\end{itemize}
\vspace{1em}

%% ── 目录 ──
\small
\tableofcontents
\normalsize
\vspace{1em}

%% ══════════════════════════════════════════════════════════════════
%%  第一模块
%% ══════════════════════════════════════════════════════════════════
\section{第一模块 \quad <模块名>}

\subsection{第1章 \quad <章名>}

\knowtitle{1.1 概念引入}

\noindent <本节要解决什么问题？为什么重要？>

\knowtitle{1.2 核心原理/公式}

\formula{$\displaystyle <核心内容>$}

\noindent\infobox{<提示>}

\warning{<常见误区>}

\knowtitle{1.3 方法技巧}

\noindent\textcolor{emphcolor}{\textbf{要点：}}
\begin{itemize}
  \item <技巧1>
  \item <技巧2>
\end{itemize}

\knowtitle{1.4 典型示例}

\begin{exampenv}{<标题>}
<示例展示>
\end{exampenv}

%% ── 流程图/状态图示例（箭头标签防重叠） ──
%% 使用 alabel 样式避免标签与目标重叠：
%%
%%   \begin{center}
%%   \begin{tikzpicture}[node distance=2.5cm]
%%     \node (A) {起点};
%%     \node (B) [right of=A] {终点};
%%     \draw[->] (A) -- node[alabel] {条件} (B);
%%   \end{tikzpicture}
%%   \end{center}
%%
%% 更多样式：alabel above（箭头上方偏左）、alabel below（箭头下方偏左）
%% \draw[->] (A) -- node[alabel above] {条件} (B);
%% \draw[->] (A) -- node[alabel below] {条件} (B);

\knowtitle{1.5 真题/实战}

\begin{exampenv}{<来源+年份+描述>}
<题目和解答>
\end{exampenv}

\knowtitle{1.6 巩固练习}

\begin{pracenv}{<标题>}
<练习>
\end{pracenv}

\knowtitle{1.7 课后作业}

\begin{hwenv}{<标题>}{ans:ch1}
<题目>
\end{hwenv}

\sep

%% ══════════════════════════════════════════════════════════════════
%%  更多章节...（复制上面的结构，注意 ans:chX 标签递推）
%% ══════════════════════════════════════════════════════════════════

%% ══════════════════════════════════════════════════════════════════
%%  专题/题型专项总结
%% ══════════════════════════════════════════════════════════════════
\section{专题/题型专项总结}

\textcolor{gray}{\small [版本 <版本号> · 最后修订 <日期>]}

~\\[0.3em]

\knowtitle{专题一：<名称>}

\noindent\textbf{方法总结}：<说明>

\formula{<核心公式/原理>}

\textcolor{supercolor}{\small \textbf{关联：}} \textcolor{supercolor}{\small <来源>}

\noindent\textbf{典型示例}：

\begin{exampenv}{<标题>}
<展示>
\end{exampenv}

\noindent\textbf{配套练习}：

\begin{pracenv}{<标题>}
<练习>
\end{pracenv}

\vspace{0.5em}

%% ── 更多专题... ──

\sep

%% ── 速查表 ──
\subsection{核心速查表}

% 按模块分子节
\subsubsection*{模块一：<名称>}

\begin{center}
{\footnotesize
\begin{tabular}{@{}ll@{}}
\toprule
\textbf{条目} & \textbf{内容} \\
\midrule
<条目A> & $\displaystyle <内容>$ \\
<条目B> & $\displaystyle <内容>$ \\
\bottomrule
\end{tabular}
}
\end{center}

\vspace{0.5em}

%% ── 作业答案 ──
\subsection{作业答案与提示}

\subsubsection*{第1章 <章名>} \label{ans:ch1}
\begin{enumerate}[label=\textbf{1.\arabic*}]
  \item <答案> \quad$\Big($说明\Big)$
\end{enumerate}

%% ── 拓展阅读/参考文献 ──
\subsection*{拓展阅读与参考文献}

\vspace{0.2em}
\noindent
\hangindent=1.5em
\hangafter=1
\textbf{[1]} <作者>.<标题>. <出版社>，<年份>. \\
\hspace*{1.5em}\textcolor{gray}{\small <备注>}

\vspace{0.5em}
\noindent\textcolor{gray}{\small <致谢> \quad <作者> \quad 最后修订 <日期>}

%% ════════════════════════════════════════════
%%  结尾
%% ════════════════════════════════════════════
\vspace{1em}
\begin{center}
  {\Huge\bfseries\color{titlecolor} \textemdash\textemdash 保持学习！\textemdash\textemdash}
\end{center}

\end{document}
